Let $(\xi_k)_{k\geq0}$ be random variables on one probability space,
let $(c_k)$ be deterministic complex weights, and consider the prefixes
\[
 f_N(z)=\sum_{k=0}^N\xi_kc_kz^k.
\]
Note that every degree uses the same coefficient sequence. For independent
centered coefficients with unit variance, the normalization is the
square root of $\sum_{k=0}^N|c_k|^2r^{2k}$.

\begin{proposition}[Polynomially weighted logarithmic variance]
\label{prop:weighted-log-variance-profile}
Fix $\tau>-1/2$, put $c_k=k^\tau$ for $k\geq1$,
and choose a fixed finite $c_0\ne0$. For the normalization
\[
 \sigma_{N,\tau}(r)^2=|c_0|^2+\sum_{k=1}^N k^{2\tau}r^{2k},
\]
we have
\[
 \log\sigma_{N,\tau}(1+x/N)-(\tau+1/2)\log N
 \longrightarrow F_\tau(x):=\frac12\log\int_0^1t^{2\tau}e^{2xt}\dd t
\]
uniformly on compact intervals of $x$. Moreover,
\[
 F_\tau'(x)=
 \frac{\int_0^1t^{2\tau+1}e^{2xt}\dd t}
      {\int_0^1t^{2\tau}e^{2xt}\dd t},
 \qquad F_\tau'(0)=\frac{2\tau+1}{2\tau+2}.
\]
\end{proposition}

The four-secant formula in Lemma~\ref{thm:T2} applies, where
$F_{N,\tau}(x)=\log\sigma_{N,\tau}(1+x/N)-(\tau+1/2)\log N$.
With the centering $\frac12\log N$, we add $\tau\log N$ to $F_{N,\tau}$,
a term that cancels in every secant. Both the convention at $k=0$ and the
restriction $\tau>-1/2$ are needed for the integral profile.

\begin{proof}
Since every compact interval of $x$ lies in some $[-K,K]$, we fix $K>0$
and work with $|x|\leq K$. The idea is to show that
$\sigma_{N,\tau}(1+x/N)^2/N^{2\tau+1}$ converges to
$\int_0^1t^{2\tau}e^{2xt}\dd t$ uniformly for $|x|\leq K$, and then to
take logarithms. For $|x|\leq K$ and $0<t\leq1$, let
\[
 g_x(t)=t^{2\tau}e^{2xt},\qquad
 I_\tau(x)=\int_0^1g_x(t)\dd t=\int_0^1t^{2\tau}e^{2xt}\dd t.
\]
For the weighted sum, we have
\[
 S_{N,\tau}(x):=\frac1{N^{2\tau+1}}\sum_{k=1}^N k^{2\tau}e^{2xk/N}
 =\frac1N\sum_{k=1}^N(k/N)^{2\tau}e^{2xk/N}
 =\frac1N\sum_{k=1}^Ng_x(k/N).
\]

\emph{Step 1: the radial factor.}
Let $|x|\leq K$ and $N\geq2K$. Multiplying the comparison
\eqref{eq:profile-comparison} by $k^{2\tau}e^{2xk/N}/N^{2\tau+1}>0$ and
summing over $1\leq k\leq N$, we get, by the definition of $S_{N,\tau}$,
\[
 e^{-2K^2/N}S_{N,\tau}(x)
 \leq\frac1{N^{2\tau+1}}\sum_{k=1}^N k^{2\tau}(1+x/N)^{2k}
 \leq e^{2K^2/N}S_{N,\tau}(x).
\]

\emph{Step 2: the Riemann sums for $\tau\geq0$.}
When $\tau\geq0$, the integrand is continuous, since the map
$(x,t)\mapsto g_x(t)$ extends continuously to the compact set
$[-K,K]\times[0,1]$, and hence it is uniformly continuous there.
Let $\eta(\delta)$ be the supremum of $|g_x(t)-g_x(t')|$ over
$|x|\leq K$ and $t,t'\in[0,1]$ with $|t-t'|\leq\delta$, so that
$\eta(\delta)\to0$ as $\delta\to0$. For $|x|\leq K$ we have
\begin{align*}
 |S_{N,\tau}(x)-I_\tau(x)|
 &=\left|\sum_{k=1}^N\frac1Ng_x(k/N)
    -\sum_{k=1}^N\int_{(k-1)/N}^{k/N}g_x(t)\dd t\right|\\
 &=\left|\sum_{k=1}^N\int_{(k-1)/N}^{k/N}
    \bigl(g_x(k/N)-g_x(t)\bigr)\dd t\right|\\
 &\leq\sum_{k=1}^N\int_{(k-1)/N}^{k/N}|g_x(k/N)-g_x(t)|\dd t
 \leq\eta(1/N),
\end{align*}
where the last step uses $|k/N-t|\leq1/N$ on the $k$-th interval, whose
lengths add up to one.

\emph{Step 3: the Riemann sums for $-1/2<\tau<0$.}
When $-1/2<\tau<0$, we fix $0<\varepsilon<1$ and split the sum at
$k_\varepsilon=\lfloor\varepsilon N\rfloor$, where $N\geq2/\varepsilon$,
so that $2\leq k_\varepsilon<N$. Then
\begin{align*}
 S_{N,\tau}(x)-I_\tau(x)
 ={}&\left[\frac1N\sum_{k=1}^{k_\varepsilon}g_x(k/N)
      -\int_0^{k_\varepsilon/N}g_x(t)\dd t\right]\\
 &+\left[\frac1N\sum_{k=k_\varepsilon+1}^Ng_x(k/N)
      -\int_{k_\varepsilon/N}^1g_x(t)\dd t\right].
\end{align*}
Since $t\mapsto t^{2\tau}$ is decreasing, we have
$\frac1N(k/N)^{2\tau}\leq\int_{(k-1)/N}^{k/N}t^{2\tau}\dd t$ for
$k\geq2$, and the term $k=1$ is $\frac1N(1/N)^{2\tau}=N^{-(2\tau+1)}$.
Also $e^{2xt}\leq e^{2K}$ for $|x|\leq K$ and $0\leq t\leq1$, and
$k_\varepsilon/N\leq\varepsilon$. Hence the normalized initial part
satisfies
\begin{align*}
 0\leq\frac1N\sum_{k=1}^{k_\varepsilon}g_x(k/N)
 &\leq\frac{e^{2K}}N\sum_{k=1}^{k_\varepsilon}(k/N)^{2\tau}
 =e^{2K}\left(N^{-(2\tau+1)}
    +\sum_{k=2}^{k_\varepsilon}\frac1N(k/N)^{2\tau}\right)\\
 &\leq e^{2K}\left(N^{-(2\tau+1)}
    +\int_{1/N}^{k_\varepsilon/N}t^{2\tau}\dd t\right)
 \leq e^{2K}\left(\frac{\varepsilon^{2\tau+1}}{2\tau+1}
                  +N^{-(2\tau+1)}\right),
\end{align*}
where the intervals $[(k-1)/N,k/N]$ with $2\leq k\leq k_\varepsilon$ make
up $[1/N,k_\varepsilon/N]$.
By the same two facts, the corresponding integral satisfies
\[
 0\leq\int_0^{k_\varepsilon/N}g_x(t)\dd t
 \leq e^{2K}\int_0^\varepsilon t^{2\tau}\dd t
 =e^{2K}\frac{\varepsilon^{2\tau+1}}{2\tau+1}.
\]
Since both initial parts are nonnegative, the first bracket is at most
$e^{2K}(\varepsilon^{2\tau+1}/(2\tau+1)+N^{-(2\tau+1)})$ in absolute value.
For the remaining Riemann sum, every interval $[(k-1)/N,k/N]$ with
$k\geq k_\varepsilon+1$ lies in $[\varepsilon/2,1]$, since
$k_\varepsilon/N\geq\varepsilon-1/N\geq\varepsilon/2$. On the compact set
$[-K,K]\times[\varepsilon/2,1]$ the map $(x,t)\mapsto g_x(t)$ is
uniformly continuous. Let $\eta_\varepsilon(\delta)$ be the
supremum of $|g_x(t)-g_x(t')|$ over $|x|\leq K$ and
$t,t'\in[\varepsilon/2,1]$ with $|t-t'|\leq\delta$, so that
$\eta_\varepsilon(\delta)\to0$ as $\delta\to0$. Then the second bracket
satisfies
\begin{align*}
 &\left|\frac1N\sum_{k=k_\varepsilon+1}^Ng_x(k/N)
      -\int_{k_\varepsilon/N}^1g_x(t)\dd t\right|
 =\left|\sum_{k=k_\varepsilon+1}^N\int_{(k-1)/N}^{k/N}
    \bigl(g_x(k/N)-g_x(t)\bigr)\dd t\right|\\
 &\qquad\leq\sum_{k=k_\varepsilon+1}^N\int_{(k-1)/N}^{k/N}
    |g_x(k/N)-g_x(t)|\dd t
 \leq\eta_\varepsilon(1/N),
\end{align*}
where the last step holds as in Step~2, since $k/N$ and $t$ lie in
$[\varepsilon/2,1]$. Therefore
\[
 \sup_{|x|\leq K}|S_{N,\tau}(x)-I_\tau(x)|
 \leq e^{2K}\left(\frac{\varepsilon^{2\tau+1}}{2\tau+1}
                  +N^{-(2\tau+1)}\right)
     +\eta_\varepsilon(1/N).
\]
Since $2\tau+1>0$, the right-hand side tends to
$e^{2K}\varepsilon^{2\tau+1}/(2\tau+1)$ as $N\to\infty$, so that the limit
superior of the left side as $N\to\infty$ is at most this value for every
$0<\varepsilon<1$, and it is therefore zero.

We conclude that in both cases $S_{N,\tau}(x)\to I_\tau(x)$ uniformly for
$|x|\leq K$.

\emph{Step 4: the variance.}
Let $|x|\leq K$ and $N\geq2K$. By the definition of $\sigma_{N,\tau}$
and Step~1, we have
\begin{align*}
 \left|\frac{\sigma_{N,\tau}(1+x/N)^2}{N^{2\tau+1}}-I_\tau(x)\right|
 \leq{}&\frac{|c_0|^2}{N^{2\tau+1}}
   +\left|\frac1{N^{2\tau+1}}\sum_{k=1}^N k^{2\tau}(1+x/N)^{2k}
   -S_{N,\tau}(x)\right|\\
 &+|S_{N,\tau}(x)-I_\tau(x)|\\
 \leq{}&\frac{|c_0|^2}{N^{2\tau+1}}+(e^{2K^2/N}-1)S_{N,\tau}(x)
   +|S_{N,\tau}(x)-I_\tau(x)|\\
 \leq{}&\frac{|c_0|^2}{N^{2\tau+1}}
   +(e^{2K^2/N}-1)\left(\frac{e^{2K}}{2\tau+1}
   +|S_{N,\tau}(x)-I_\tau(x)|\right)\\
 &+|S_{N,\tau}(x)-I_\tau(x)|,
\end{align*}
where the second inequality also uses $1-e^{-2K^2/N}\leq e^{2K^2/N}-1$,
which holds since $e^u+e^{-u}\geq2$, and the last uses
$I_\tau(x)\leq e^{2K}\int_0^1t^{2\tau}\dd t=e^{2K}/(2\tau+1)$.
Since $2\tau+1>0$, the term with the fixed $c_0$ also vanishes after
normalization, and by Steps~2 and~3 the other terms tend to zero
uniformly for $|x|\leq K$, since $e^{2K^2/N}-1\to0$. Thus
$\sigma_{N,\tau}(1+x/N)^2/N^{2\tau+1}\to\int_0^1t^{2\tau}e^{2xt}\dd t$
uniformly for $|x|\leq K$.

\emph{Step 5: the logarithm.}
For $|x|\leq K$, the integral is at least $e^{-2K}/(2\tau+1)>0$, since
$I_\tau(x)\geq e^{-2K}\int_0^1t^{2\tau}\dd t$. By Step~1 we have
\begin{align*}
 \frac{\sigma_{N,\tau}(1+x/N)^2}{N^{2\tau+1}}
 &\geq\frac1{N^{2\tau+1}}\sum_{k=1}^Nk^{2\tau}(1+x/N)^{2k}
 \geq e^{-2K^2/N}S_{N,\tau}(x)\\
 &\geq e^{-2K^2/N}\bigl(I_\tau(x)-|S_{N,\tau}(x)-I_\tau(x)|\bigr)\\
 &\geq e^{-2K^2/N}\left(\frac{e^{-2K}}{2\tau+1}
      -|S_{N,\tau}(x)-I_\tau(x)|\right),
\end{align*}
where the first inequality drops $|c_0|^2/N^{2\tau+1}\geq0$.
Also $I_\tau(x)$ is at least the same lower bound, since
$e^{-2K^2/N}\leq1$. Since
$|\log a-\log b|\leq|a-b|/\min(a,b)$ for $a,b>0$, we get, for all $N$
so large that this lower bound is positive for $|x|\leq K$,
\begin{align*}
 &\left|\log\sigma_{N,\tau}(1+x/N)-(\tau+1/2)\log N-F_\tau(x)\right|\\
 &\qquad=\frac12\left|\log\frac{\sigma_{N,\tau}(1+x/N)^2}{N^{2\tau+1}}
    -\log I_\tau(x)\right|\\
 &\qquad\leq\frac{\bigl|\sigma_{N,\tau}(1+x/N)^2/N^{2\tau+1}-I_\tau(x)\bigr|}
   {2e^{-2K^2/N}\bigl(e^{-2K}/(2\tau+1)-|S_{N,\tau}(x)-I_\tau(x)|\bigr)},
\end{align*}
since $F_\tau=\frac12\log I_\tau$ and both arguments of the logarithms are
at least the lower bound above.
By Steps~2 to~4, the numerator tends to zero and the denominator tends
to $2e^{-2K}/(2\tau+1)$, both uniformly for $|x|\leq K$. Hence the left
side tends to zero uniformly for $|x|\le K$, which proves the profile
limit of the proposition.

By differentiating the positive integral, we get
\[
 \frac{\mathrm d}{\mathrm dx}\left(\frac12\log
                \int_0^1t^{2\tau}e^{2xt}\dd t\right)
 =\frac12\cdot\frac{\int_0^12t\cdot t^{2\tau}e^{2xt}\dd t}
        {\int_0^1t^{2\tau}e^{2xt}\dd t}
 =\frac{\int_0^1t^{2\tau+1}e^{2xt}\dd t}
        {\int_0^1t^{2\tau}e^{2xt}\dd t},
\]
where differentiation under the integral sign is allowed since, for
$|x|\leq K$, the $x$-derivative $2t^{2\tau+1}e^{2xt}$ of the integrand is
at most $2e^{2K}t^{2\tau}$, which is integrable on $[0,1]$ since
$2\tau>-1$. At zero this quotient is
$\int_0^1t^{2\tau+1}\dd t/\int_0^1t^{2\tau}\dd t=(2\tau+1)/(2\tau+2)$.
\end{proof}

\begin{remark}\label{rem:weighted-divergence}
For $\tau\leq-1/2$, the integral defining $F_\tau$ diverges at zero.
The finite-degree radial secant estimate still applies with the exact
normalization. To deduce an almost-sure root law from a variance profile,
one must also establish the occupation and logarithmic hypotheses and
the degree-block estimates.
\end{remark}
The integral diverges since its integrand is at least
$e^{-2|x|}t^{2\tau}$ and $2\tau\leq-1$. The finite-degree radial secant
estimate of the remark is \eqref{eq:T2-exact} of Lemma~\ref{thm:T2},
which allows any normalization $\sigma_N(r)>0$, and in particular
$\sigma_{N,\tau}$.

